%%%%%%%%%%%%%%%%%%%%%%%%%%%%%%%%%%%%%%%%%
% FRI Data Science_report LaTeX Template
% Project report for the NLP course
%%%%%%%%%%%%%%%%%%%%%%%%%%%%%%%%%%%%%%%%%

\documentclass[fleqn,moreauthors,10pt]{ds_report}
\usepackage[english]{babel}
\usepackage[T1]{fontenc}
\usepackage{tabularx}
\usepackage{array}

\graphicspath{{fig/}}

\JournalInfo{FRI Natural Language Processing Course 2025/2026}
\Archive{Project report}
\PaperTitle{Zakonodajko: Retrieval-Augmented Question Answering over Slovenian Tax Legislation}
\Authors{ByteBanda project team}
\affiliation{\textit{Advisor: Slavko \v{Z}itnik}}
\Keywords{retrieval-augmented generation, Slovenian legislation, tax law, FAISS, sentence embeddings}
\newcommand{\keywordname}{Keywords}
\newcommand{\code}[1]{\texttt{#1}}
\newcolumntype{Y}{>{\raggedright\arraybackslash}X}
\setlength{\textfloatsep}{6pt plus 1pt minus 2pt}
\setlength{\floatsep}{6pt plus 1pt minus 2pt}
\setlength{\intextsep}{6pt plus 1pt minus 2pt}
\hypersetup{
    pdftitle={Zakonodajko: Retrieval-Augmented Question Answering over Slovenian Tax Legislation},
    pdfauthor={ByteBanda project team}
}

\Abstract{
Zakonodajko is a Retrieval-Augmented Generation (RAG) assistant for Slovenian tax legislation. It extracts local PISRS HTML documents, chunks them, embeds the chunks with a multilingual sentence-transformer, indexes them in FAISS, retrieves grounded context, and prompts a local instruction-tuned language model. We compare a fixed dense baseline with a Version 2 pipeline that adds article-aware legal chunking and hybrid reranking. The corrected aggregate results from the existing report show that Version 2 improves source hit@3 from 40.0\% to 90.0\% for the full 30-question run, while the strict 20-question prompt run reaches 95.0\% source hit@3, 40.0\% article hit@3, and 0.65/2 manual correctness.
}

\begin{document}

\flushbottom
\maketitle
\thispagestyle{empty}

\section*{Introduction}

Legal question answering requires exact source grounding: a wrong article, exception, threshold, or act can change the answer. Our project, \textit{Zakonodajko}, addresses this by using RAG over Slovenian tax legislation instead of relying on model memory alone. The collection contains local PISRS versions of ZDDV-1, ZDoh-2, ZDDPO-2, ZDavP-2, and related implementing rulebooks. We kept the system retrieval-first because the main bottleneck was retrieval precision, not model capacity.

\section*{Methods}

The baseline system splits documents into overlapping 1,200-character windows with 200 characters of overlap, embeds chunks with multilingual MiniLM, normalizes vectors, and searches a FAISS \code{IndexFlatIP} index with \code{top\_k=3}. Version 2 keeps dense retrieval as the candidate generator but changes two parts of the pipeline. First, article-aware chunking detects headings such as \code{3. \v{c}len} and \code{113.a \v{c}len}, then stores source filename, law identifier, document role, article number, article title, and section position as metadata. Second, hybrid reranking scores the top \code{candidate\_k=100} dense candidates as
\[
s(q,c)=s_{\mathrm{dense}}(q,c)+0.25s_{\mathrm{lexical}}(q,c)+s_{\mathrm{law}}(q,c),
\]
where the lexical term is normalized token overlap over chunk text plus metadata, and the law boost is 0.20 when the question names or strongly implies the same legal act as the chunk.

The generation step renders the retrieved chunks with source, chunk ID, score, law ID, article number, and title. We evaluated the baseline prompt, which asks the model to answer only from context and cite sources, and a strict prompt, which tells the model to prefer the chunk whose act and article directly answer the question and to say the answer is not found when the exact rule is missing. Retrieval is evaluated with source hit@3, article hit@3, phrase hit@3, and manual context relevance. Generation is scored manually from 0 to 2 for faithfulness and correctness.

\begin{table}[hbt]
\caption{Legal documents used in the RAG evaluation.}
\centering
\small
\begin{tabularx}{\linewidth}{lY}
\toprule
ID & Document \\
\midrule
ZDDV-1 & Value Added Tax Act and rulebook \\
ZDoh-2 & Personal Income Tax Act \\
ZDDPO-2 & Corporate Income Tax Act and rulebook \\
ZDavP-2 & Tax Procedure Act and rulebook \\
\bottomrule
\end{tabularx}
\label{tab:data}
\end{table}

\section*{Results}

Table~\ref{tab:retrieval} uses the aggregate values from Table 2 of the existing \code{report.pdf}. The corrected values show that the full Version 2 legal-hybrid retrieval run improves source hit@3 to 0.90, article hit@3 to 0.30, and phrase hit@3 to 0.35 over 30 questions. The strict prompt was available for 20 questions and reaches the best source and article scores in the table: 0.95 source hit@3 and 0.40 article hit@3.

\begin{table}[hbt]
\caption{Automatic evaluation metrics from the existing report results. Src/Art/Phr are @3 metrics; Ctx/Fth/Cor use 0--2 scores.}
\centering
\scriptsize
\setlength{\tabcolsep}{2.5pt}
\begin{tabular}{lrrrrrrr}
\toprule
Run & n & Src & Art & Phr & Ctx & Fth & Cor \\
\midrule
V1 old & 30 & 0.40 & 0.10 & 0.15 & 0.10 & 1.00 & 0.25 \\
V2 fixed & 30 & 0.45 & 0.15 & 0.20 & 0.40 & 1.05 & 0.30 \\
V2 hybrid & 30 & 0.90 & 0.30 & 0.35 & 1.30 & 1.15 & 0.40 \\
V2 base & 30 & 0.90 & 0.35 & 0.35 & 1.30 & 1.45 & 0.65 \\
V2 strict & 20 & 0.95 & 0.40 & 0.35 & 1.30 & 1.05 & 0.80 \\
\bottomrule
\end{tabular}
\label{tab:retrieval}
\end{table}

The most important improvement is source-level retrieval: the legal-hybrid retriever usually selects the correct act, while the fixed dense runs often retrieve semantically related material from rulebooks or another law. Article retrieval remains the main limitation. Even in the best strict-prompt run, article hit@3 is 0.40, so the system often finds the right law but not the exact definitional or procedural article.

\begin{table}[hbt]
\caption{Manual generation assessment from Table 3 of the existing report. Scores use a 0--2 scale over 20 questions.}
\centering
\small
\begin{tabular}{lrr}
\toprule
Metric & Baseline prompt & Strict prompt \\
\midrule
Context relevance & 1.20 & 1.25 \\
Faithfulness & 1.45 & 1.55 \\
Correctness & 0.60 & 0.65 \\
\bottomrule
\end{tabular}
\label{tab:manual}
\end{table}

Table~\ref{tab:manual} shows that strict prompting gives only a small correctness gain. This is expected because both prompt variants depend on the same retrieved context. When the correct article is present, both prompts can answer well; when retrieval returns only related material, the model either gives a faithful but legally wrong answer or refuses because the exact rule is missing. The common error patterns are: correct law but wrong article, rulebook provisions instead of the statute definition, cross-law confusion between ZDoh-2 and ZDDPO-2, and faithful answers to irrelevant retrieved passages.

\section*{Discussion and Conclusion}

This report should be read as a midpoint checkpoint, not as the final evaluation of the project. The current experiments exposed an important evaluation and retrieval bias: many generated evaluation questions explicitly contain law names such as ZDDV-1, ZDoh-2, ZDDPO-2, or ZDavP-2, and the retriever uses a matching law-name signal as an additional score. This makes source correctness much higher, but it does not necessarily prove that the retriever understands the legal question or can find the exact article without this clue. Exact article matching remains unreliable, and the existing evaluation data is not complete enough to measure legal answer quality robustly.

The next stage will therefore focus on reducing this bias and building a complete human-made evaluation dataset with validated questions, exact expected sources and articles, key legal phrases, and reference answers. This should make evaluation independent of old chunk IDs and less dependent on questions that reveal the correct law in advance. In parallel, we plan to extend the RAG pipeline with additional components such as tool use, query rewriting, answer verification, or citation checking. Finally, the generation experiments will be repeated with the GAMS LLM instead of Mistral, because GAMS is a Slovenian model and should be better suited for producing grammatically and semantically correct Slovene answers.

\end{document}
